% Section 8 -- Comparison to guard models. Framing rule 6: the case is cost profile, NOT accuracy.
% The guard model wins on accuracy and the section says so plainly, in its own words, in the first paragraph.
\section{Comparison to Guard Models}
\label{sec:guard}

The deployed alternative to an activation probe is a guard model: a separate classifier, itself a language model,
that reads the prompt and returns a safety verdict. Llama Guard~3 \cite{inan2023llama} is the common open reference,
and ShieldGemma \cite{zeng2024shieldgemma} and WildGuard \cite{han2024wildguard} occupy the same slot. The comparison
below runs both systems on identical inputs through their own intended interfaces.

\textbf{The guard model is more accurate, on every metric, and the probe's case is not accuracy.}
Table~\ref{tab:guard} gives the numbers: Llama Guard~3 reaches $93.5\,\%$ accuracy against the probe's $81.7\,\%$,
with a false-positive rate of $8.9\,\%$ against $26.9\,\%$ and F1 of $0.932$ against $0.822$. We state this first
because the comparison is easy to present the other way round by choosing a favourable prompt set, and because the
probe's actual advantage does not depend on winning it.

% GENERATED by paper/tools/gen_tables.py -- do not edit by hand.
% sources: results/phase1/llama_guard_comparison.json
\begin{table*}[!t]
\caption{Probe versus guard model on the same $650$ prompts ($300$ HarmBench, $250$ XSTest safe, $100$
JailbreakBench benign), each through its own intended interface. The guard model is more accurate on every
metric; the probe's case is cost, not accuracy.}
\label{tab:guard}
\centering
\small
\begin{tabular}{lccccccc}
\toprule
System & Acc.\ & TPR & FPR & Prec. & F1 & latency (ms) & p99 \\
\midrule
Probe on Llama-3.2-3B & 81.7 & 91.7 & 26.9 & 0.745 & 0.822 & 16.2 $\pm$ 1.2 & 21.1 \\
Llama Guard 3 (8B) & 93.5 & 96.3 & 8.9 & 0.903 & 0.932 & 64.7 $\pm$ 21.5 & 107.0 \\
\bottomrule
\end{tabular}

\footnotesize Mean-latency ratio $4.0\times$ on identical hardware (A100-40GB), over 3 runs. All $94$ of the probe's false positives are JailbreakBench benign prompts.
\end{table*}

The evaluation is $650$ prompts---$300$ HarmBench behaviours, $250$ XSTest safe prompts, $100$ JailbreakBench benign
prompts---with the guard model prompted through its own chat template including its hazard taxonomy, and the probe
read at its own layer and threshold. Both systems see exactly the same text.

The shape of the probe's errors is worth more than the aggregate. All $94$ of its false positives are JailbreakBench
benign prompts, and none are XSTest: the failure is entirely on on-topic benign requests, which is the topic-confound
signature of Section~\ref{sec:learn} appearing in a system-level metric. Llama Guard's false positives are spread
more evenly ($3.2\,\%$ on XSTest, $23.0\,\%$ on JailbreakBench), so it is not merely a better-calibrated version of
the same detector; it is discriminating on a partly different basis.

\subsection{Where the cost actually sits}
\label{sec:cost}

% GENERATED by paper/tools/gen_tables.py -- do not edit by hand.
% sources: results/phase1/latency/*.json
\begin{table}[!t]
\caption{Where the cost sits. The probe itself is a dot product; the forward pass it rides on is the expense,
and that pass is already being computed to serve the request. Batch size $1$, greedy, hook attached, $50$ timed
generations after $5$ warm-ups, on an A100-40GB.}
\label{tab:latency}
\centering
\setlength{\tabcolsep}{4pt}
\footnotesize
\begin{tabular}{lccc}
\toprule
Model & fwd.\ pass (ms) & p99 & probe ($\mu$s) \\
\midrule
Gemma-2-2B & 31.5 $\pm$ 0.2 & 32.1 & 1.1 \\
Gemma-2-9B & 50.0 $\pm$ 0.2 & 50.3 & 1.3 \\
Gemma-3-1B & 41.5 $\pm$ 0.2 & 42.2 & 1.0 \\
Gemma-3-4B & 61.2 $\pm$ 1.6 & 68.9 & 1.3 \\
Llama-3.1-8B & 17.7 $\pm$ 0.2 & 18.1 & 1.5 \\
Llama-3.2-1B & 10.1 $\pm$ 0.1 & 10.5 & 1.2 \\
Llama-3.2-3B & 15.9 $\pm$ 0.2 & 16.5 & 1.4 \\
\bottomrule
\end{tabular}
\end{table}

The probe's case is the cost profile, and Table~\ref{tab:latency} shows why it is not the ratio in
Table~\ref{tab:guard} that matters. The probe itself is a dot product between a $3{,}584$-dimensional activation and a
direction vector: measured at $1.2$--$1.5$ microseconds, four orders of magnitude below the forward pass it rides on.
And that forward pass is not an additional cost, because it is the pass already being run to serve the request. In a
deployment that is generating a response anyway, the marginal cost of the probe is the dot product and a hook.

A guard model has a different structure of cost: it is a second model, of comparable size to the one being guarded,
run as an extra inference. Here that is $64.7$~ms mean and $107.0$~ms at the ninety-ninth percentile against the
probe path's $16.2$~ms and $21.1$~ms, a mean ratio of $4.0\times$; but the ratio is an artifact of which host model
the probe rides on, and would change with a different one. The durable statement is architectural rather than
numerical: a probe adds microseconds to an inference that is already happening, whereas a guard model adds an
inference. Memory follows the same pattern---an $8$B guard model must be resident alongside the served model, while a
direction vector is a few kilobytes.

\textbf{Hardware caveat.} These are A100-40GB numbers on a current inference stack. Earlier figures for this codebase
were collected on an L4, where the absolute latencies and the ratio both differ; the two are not comparable, and we
report only the measurements made here. Latency ratios between a probe and a guard model should be read as a property
of a particular pairing on particular hardware, not as a constant of the method.

\subsection{What the comparison supports}
\label{sec:guard_claims}

An activation probe is not a cheaper way to get guard-model accuracy, and on this evidence it is not close. What it
offers is a detector whose marginal cost is negligible on an inference that is already being performed, which makes
it a candidate for a first-stage filter or an always-on monitor rather than a replacement for a guard model. Whether
that is worth $12$ accuracy points, and a false-positive rate concentrated precisely on benign requests that share a
topic with harmful ones, is a deployment question rather than a research one---but it should be asked with the honest
numbers in Table~\ref{tab:guard} rather than with a comparison tuned to favour the probe.
